\documentclass[11pt,letterpaper]{article}

\usepackage[margin=0.72in,headheight=15pt]{geometry}
\usepackage[T1]{fontenc}
\usepackage{lmodern}
\usepackage{microtype}
\usepackage{xcolor}
\usepackage{amsmath,amssymb}
\usepackage{enumitem}
\usepackage{tabularx,booktabs,longtable,array}
\usepackage{fancyhdr}
\usepackage{hyperref}
\usepackage{listings}

\definecolor{ink}{HTML}{17324D}
\definecolor{blue}{HTML}{126B8A}
\definecolor{pale}{HTML}{EAF4F7}
\definecolor{gold}{HTML}{D89614}
\definecolor{gray}{HTML}{56616C}
\definecolor{line}{HTML}{C9D4DC}
\hypersetup{colorlinks=true,linkcolor=blue,urlcolor=blue,citecolor=blue,pdftitle={48-Hour BlueQubit Challenge Practice Handout},pdfauthor={Team practice handout}}
\pagestyle{fancy}
\fancyhf{}
\fancyhead[L]{\small\textcolor{ink}{DUQUANTUM / BLUEQUBIT}}
\fancyhead[R]{\small\textcolor{gray}{48-hour practice handout}}
\fancyfoot[C]{\small\textcolor{gray}{\thepage}}
\setlength{\parindent}{0pt}
\setlength{\parskip}{5pt}
\setlist[itemize]{leftmargin=1.35em,itemsep=2pt,topsep=3pt}
\setlist[enumerate]{leftmargin=1.55em,itemsep=4pt,topsep=3pt}
\renewcommand{\arraystretch}{1.18}

\newcommand{\taskbox}[1]{\vspace{2pt}\noindent\colorbox{pale}{\parbox{\dimexpr\linewidth-2\fboxsep}{\textcolor{ink}{\textbf{Deliverable:} #1}}}\vspace{3pt}}
\newcommand{\pass}[1]{\textbf{Pass check:} #1}
\newcommand{\source}[2]{\textbf{Source:} \href{#1}{#2}}
\newcommand{\checkline}{\(\square\) }
\newcommand{\role}[2]{\textbf{#1}\hfill #2\\}

\begin{document}

\begin{center}
{\Large\textcolor{blue}{48-HOUR PRACTICE PLAYBOOK}}\\[4pt]
{\huge\bfseries\textcolor{ink}{Estimate the Observables}}\\[4pt]
{\large A source-grounded team handout for the BlueQubit challenge at DuQuantum}\\[8pt]
{\small Duke University \textbullet\ October 24--25, 2026 \textbullet\ Prepare in two focused days}
\end{center}
\vspace{4pt}\hrule\vspace{7pt}

\textbf{Use this as a work plan, not a quantum textbook.} The fastest route to a strong challenge submission is to build a checked estimation pipeline, learn the simulator controls BlueQubit exposes, and benchmark the exact score under a fixed compute budget. Your linear algebra and VQC background is enough for the math below. Spend practice time writing code, checking outputs, and making decisions from measured results.

\begin{quote}\small
\textbf{What is publicly known (verify again at kickoff).} BlueQubit's public announcement describes about 10 hidden circuits and about 50 observable expectation values, with runs made through BlueQubit by circuit ID. It says the team can choose methods and compute spend, and that scoring depends on correlation with exact values / signal-to-noise. DuQuantum's site lists a 24-hour event, October 24--25, and teams of up to four. BlueQubit's public post describes an online option and teams up to five. Confirm which rules apply to your entry, and copy the exact scoring formula before optimizing. These are current public descriptions, not a substitute for the live challenge rules.\
\source{https://qa.linkedin.com/company/bluequbit}{BlueQubit public announcement} \quad
\source{https://duquantum.org/}{DuQuantum event and FAQ}
\end{quote}

\textbf{Strategic scope.} This handout targets the observable-estimation challenge. Do not spend these two days implementing Shor or Grover. BlueQubit's separate peaked-circuit / Bitcoin challenge is a different task; its page is marked completed.\footnote{\url{https://app.bluequbit.io/hackathons/GFgHTGbTylwmMsCp}} The most relevant skills here are Pauli expectation values, measurement, simulator selection, approximation controls, score calculation, cost tracking, and fast verification.

\section*{1. Two-day finish line}
By the end of the sprint, the team should have these concrete artifacts:
\begin{itemize}
  \item A runnable local baseline that estimates Pauli expectations and labels its qubit / bit ordering.
  \item A BlueQubit smoke-test script for a small Qiskit circuit and a short table comparing at least two simulator approaches.
  \item A score calculator that implements the challenge's actual metric after it is published.
  \item A benchmark log recording circuit, observable, backend, settings, estimate, runtime, and cost or credit usage.
  \item A one-page first-hour runbook, named team roles, and a concise method explanation ready for submission.
\end{itemize}

\textbf{Common minimum for every teammate:} Each person can explain \(\langle O\rangle=\langle\psi|O|\psi\rangle\), estimate a Pauli expectation from counts, identify the right basis rotation for X and Y, read the benchmark table, and explain the team score. Assign specialist implementation tasks, but do not leave the core pipeline understood by only one person.

\subsection*{Pocket reference: what the estimator computes}
For a state \(|\psi\rangle=U|0\ldots0\rangle\), the target is \(\langle O\rangle=\langle\psi|O|\psi\rangle\). For a Pauli string measured after the appropriate basis rotations, each shot gives a sign. If the measured bits on the non-identity positions of \(P\) have parity \(p_s\), that shot contributes \((-1)^{p_s}\), and
\[
  \widehat{\langle P\rangle}=\frac{1}{N}\sum_{s=1}^{N}(-1)^{p_s}.
\]
For an observable written as a Pauli sum \(O=\sum_k c_kP_k\), combine the term estimates as \(\widehat{\langle O\rangle}=\sum_k c_k\widehat{\langle P_k\rangle}\). A computational-basis readout directly measures Z terms; apply \(H\) before readout for X, and \(S^\dagger\) then \(H\) for Y. For a single Pauli term with outcomes \(\pm1\), the sampling standard error is approximately \(\sqrt{(1-\mu^2)/N}\), with true mean \(\mu\). This is sampling error; an approximate simulator can add a separate error source.\
\source{https://quantum.cloud.ibm.com/docs/en/guides/specify-observables-pauli}{IBM Quantum: Pauli observables and basis changes} \quad
\source{https://learning.quantum.ibm.com/course/variational-algorithm-design/cost-functions}{IBM Quantum Learning: expectation values}

\section*{2. The 48-hour schedule}
The schedule assumes about 7--8 focused practice hours per day, with breaks and sleep. If time gets tight, complete the items marked \textbf{CORE} and drop the optional reading before dropping verification.

\subsection*{Day 1: build and verify the estimator}
\begin{tabularx}{\linewidth}{@{}p{0.91in}p{1.43in}X@{}}
\toprule
\textbf{Block} & \textbf{Work} & \textbf{Exit condition}\\
\midrule
09:00--09:30 & Setup \textbf{CORE} & Everyone can run the same small local Qiskit circuit; token is stored outside source control.\\
09:30--11:00 & Practice Problems 1--3 & Statevector and measurement basics are checked on published exercises.\\
11:15--12:15 & Practice Problems 4--5 & A teammate can estimate X and Y from Z-basis readout after rotations.\\
13:00--14:00 & Practice Problem 6 & The team can compute a Hamiltonian expectation from Pauli terms.\\
14:00--15:30 & Practice Problems 7--8 & The team has tried the BlueQubit expectation API and compared methods on the same circuit.\\
15:45--16:30 & Result review & The team records every mismatch, qubit-order assumption, and simulator setting.\\
16:30--17:00 & Teach-back & Each teammate explains one exercise and one common implementation failure.\\
\bottomrule
\end{tabularx}

\subsection*{Day 2: benchmark for the score and rehearse the event}
\begin{tabularx}{\linewidth}{@{}p{0.91in}p{1.43in}X@{}}
\toprule
\textbf{Block} & \textbf{Work} & \textbf{Exit condition}\\
\midrule
09:00--10:00 & Finish Practice Problem 8 & Pauli-path and MPS calls use the same circuit and Pauli sum; results and settings are logged.\\
10:00--11:30 & Practice Problem 9 & A truncation / bond-dimension sweep is plotted or tabulated against runtime and result changes.\\
11:45--12:30 & Practice Problem 10 & The score calculator returns the expected metric on the source's own example or on recorded reference outputs.\\
13:15--14:30 & Practice Problem 11 & One teammate runs a mock challenge without revealing the circuit until the run begins.\\
14:30--15:15 & Agent-assisted review & Agents find issues or benchmark alternatives; a human checks every result and integrates one change at a time.\\
15:30--16:15 & First-hour runbook & Team roles, spending limit, smoke test, escalation path, and submission owner are written down.\\
16:15--16:45 & Oral drill & Each person answers the challenge questions in Section 7 without notes.\\
\bottomrule
\end{tabularx}

\section*{3. Source-based practice problems}
The problem statements below are paraphrases of exercises or runnable examples in the linked courseware and official product docs. Work from the linked page or notebook; do not substitute invented answer keys. Each item names a checkable output. Problems 1--10 are hands-on. Problem 11 is an event rehearsal using the actual challenge's published constraints.

\subsection*{Problem 1 -- Prepare and inspect three single-qubit states (CORE, 20 min)}
Complete IBM Quantum Learning, Utility-Scale Quantum Computing, Lesson 2, Exercise 1. Build circuits for \(XX|0\rangle\), \(HH|0\rangle\), and \(HZH|0\rangle\). Display each resulting statevector and Bloch sphere. Before running, predict the state after each gate sequence.\
\source{https://learning.quantum.ibm.com/course/utility-scale-quantum-computing/lesson-02-bits-gates-and-circuits}{Lesson 2: Bits, gates, and circuits}
\taskbox{Save the three circuits, printed statevectors, and Bloch sphere output. Explain why applying a gate twice can return the original state for the first two cases.}
\pass{Your saved results agree with the exercise's worked outputs; you can describe the gate order used by Qiskit.}

\subsection*{Problem 2 -- Bell and GHZ circuits on two simulator families (CORE, 30--40 min)}
Complete the BlueQubit BQ 101 tutorial section that builds and analyzes Bell and GHZ circuits. Run the tutorial's canonical examples with its statevector and MPS backends. Inspect the returned state or measurement outputs, then record which outputs agree and how the circuit creates entanglement.\
\source{https://app.bluequbit.io/tutorials}{BlueQubit tutorials, BQ 101}
\taskbox{Keep the tutorial's Bell/GHZ examples and backend outputs. Add a row for each backend, circuit, returned quantity, runtime, and any error.}
\pass{The team can explain why one-qubit marginals do not by themselves describe the Bell/GHZ correlations.}

\subsection*{Problem 3 -- Estimate the Bloch-vector X component (CORE, 25 min)}
Complete IBM Quantum Learning, Lesson 3, Exercise 1. Use the state preparation in the lesson, rotate before the computational-basis measurement, estimate \(r_x\), and compare your estimate with the value calculated from the preparation parameters.\
\source{https://learning.quantum.ibm.com/course/utility-scale-quantum-computing/lesson-03-teleportation}{Lesson 3: Teleportation, Exercise 1}
\taskbox{Save the measurement counts and the exact and sampled estimates. Explain which pre-measurement rotation maps the X measurement to the Z readout.}
\pass{Your estimator follows the lesson's method and your comparison includes the shot count.}

\subsection*{Problem 4 -- Estimate the Bloch-vector Y component (CORE, 25 min)}
Complete IBM Quantum Learning, Lesson 3, Exercise 2. Adapt the preceding tomography code to estimate \(r_y\) with the lesson's proposed basis changes. Compare with the parameter-derived reference.\
\source{https://learning.quantum.ibm.com/course/utility-scale-quantum-computing/lesson-03-teleportation}{Lesson 3: Teleportation, Exercise 2}
\taskbox{Save the circuit, counts, exact reference, sampled estimate, and the gate sequence immediately before measurement.}
\pass{The team can explain why the Y-basis rotation sequence and its order matter.}

\subsection*{Problem 5 -- Decompose and estimate a Hamiltonian (CORE, 30 min)}
Work the IBM Variational Algorithm Design, Cost Functions guided example. The source prepares \(|+\rangle\), gives \(H=\begin{pmatrix}-1&2\\2&1\end{pmatrix}=2X-Z\), and asks you to obtain its expectation by estimating the Pauli terms. Calculate the result both from the matrix and from \(2\langle X\rangle-\langle Z\rangle\).\
\source{https://learning.quantum.ibm.com/course/variational-algorithm-design/cost-functions}{Cost Functions, guided expectation-value example}
\taskbox{Write down the Pauli decomposition, the two separate term values, the weighted sum, and a code result using the Qiskit or BlueQubit path used by the lesson.}
\pass{The decomposition and result match the source example; the team can explain why X and Z use different measurement bases.}

\subsection*{Problem 6 -- Run the BlueQubit SDK quickstart (CORE, 30 min)}
Complete the BlueQubit SDK quickstart: initialize the client, submit the documented two-qubit Qiskit circuit, retrieve its statevector, and run the documented 40-qubit GHZ example on MPS with the documented bond dimension. Use a fresh test job name and record the actual runtime and returned type. Do not assume pricing or service availability from a saved copy of the docs; check the live challenge credits.\
\source{https://app.bluequbit.io/sdk-docs/index.html}{BlueQubit SDK quickstart, examples 1--6}
\taskbox{A short script or notebook that shows client initialization, a successful small run, result extraction, and the 40-qubit MPS call.}
\pass{The team knows how to submit, wait for, inspect, and cancel a job, and can locate the current service limits.}

\subsection*{Problem 7 -- Compare Pauli-path with MPS for the same observable (CORE, 45 min)}
Complete the BlueQubit SDK docs' Pauli-path example. Reproduce the documented three-qubit circuit and Pauli sum \texttt{XYZ}, \texttt{XXX}, \texttt{XII}, and \texttt{III} with the documented coefficients. Run it with the documented Pauli-path threshold, then compare its returned expectation with the same circuit and sum on \texttt{mps.cpu}.\
\source{https://app.bluequbit.io/sdk-docs/index.html}{BlueQubit SDK quickstart, example 7}
\taskbox{Save circuit source, Pauli sum, backend settings, both results, job IDs, and runtimes. Keep the exact source example intact as a reproducible baseline.}
\pass{You can demonstrate a same-input comparison and explain which result is documented as exact in this example.}

\subsection*{Problem 8 -- Reproduce a parameterized CHSH expectation sweep (CORE for one specialist, 60--90 min)}
Complete the IBM CHSH Inequality tutorial. Prepare the Bell pair, construct the parameterized measurement circuit and correlators, sweep the angle values listed in the tutorial from \(0\) to \(2\pi\), then compute and plot the CHSH quantities. The purpose here is to practice many expectation estimates and post-processing, not to study Bell tests in depth.\
\source{https://learning.quantum.ibm.com/tutorial/chsh-inequality}{CHSH Inequality tutorial}
\taskbox{Save the parameterized circuit, observable list, angle array, results table, and plot. Mark the angles where the tutorial's CHSH inequality is violated.}
\pass{The plotted values and interpretation match the tutorial. One teammate can describe how the expectation values feed the final score.}

\subsection*{Problem 9 -- Sweep approximation controls (CORE for simulator specialist, 60 min)}
Use the same documented BlueQubit Pauli-path example from Problem 7. Change only \texttt{pauli\_path\_truncation\_threshold} across multiple documented-valid settings. Compare each result to the MPS reference and record runtime. Then, on a circuit supported by the MPS example, vary \texttt{mps\_bond\_dimension} and log the corresponding result. The SDK documents that lower Pauli-path thresholds retain more terms at higher runtime cost; MPS bond dimension controls retained Schmidt coefficients.\
\source{https://app.bluequbit.io/sdk-docs/bluequbit.sdk.html}{BlueQubit SDK API: simulator options}
\taskbox{A table with one row per setting: backend, truncation/bond setting, expectation, reference difference, runtime, cost, and job identifier. No unbounded sweeps.}
\pass{You can choose a setting from measured tradeoffs and explain uncertainty. Do not claim the tradeoff is monotonic unless your data supports it.}

\subsection*{Problem 10 -- Match the challenge's score before tuning (CORE, 30 min)}
Read the live BlueQubit challenge instructions and implement the score exactly as written. Validate your code using the instructions' own examples or a small saved dataset for which you independently calculated the result. If the event score uses correlation, compare correlation with absolute error and confirm that your code handles constant arrays, sign reversals, and missing values as the official rules specify. Do not assume a score formula from a social post is the final formula.\
\source{https://app.bluequbit.io/hackathons}{BlueQubit hackathons; open the active Duke challenge}
\taskbox{A small metric function, one hand-calculated check, and a copied version of the official scoring definition with the source page and access time.}
\pass{Two teammates independently compute the same score for the saved check.}

\subsection*{Problem 11 -- Timed hidden-circuit rehearsal (CORE, 60--90 min)}
Use a BlueQubit tutorial circuit and its documented observables as the mock input. One teammate hides the circuit identity and gives the solver only the interface and available information that a challenge participant would receive. The solver must: record the input, run a cheap baseline, pick a method, check result format, calculate the metric wherever ground truth is available, and prepare a method note. This is a workflow rehearsal, not a claim that a tutorial circuit is a challenge instance.\
\source{https://qa.linkedin.com/company/bluequbit}{BlueQubit public challenge announcement}
\taskbox{A timed run log, a result artifact, a score if reference values are available, and a 5-sentence explanation of the method and limitations.}
\pass{The team produces a checked baseline before scaling up and finishes inside its fixed practice budget.}

\subsection*{Optional capstone -- Read the Pauli-Path tutorial (30--60 min)}
Read BlueQubit's Pauli-Path Simulation tutorial and identify the circuit, observable, truncation control, and reported comparison. Recreate only the tutorial's initial small example unless the team has spare time and credits for the larger experiment. Treat the 127-qubit demonstration as a reading exercise, not a required two-day implementation.\
\source{https://app.bluequbit.io/tutorials}{BlueQubit tutorials, Pauli-Path Simulation of Quantum Circuits}

\section*{4. Method selection card}
Use this as a hypothesis checklist. The data from your benchmark, not the names below, chooses the final method.
\begin{tabularx}{\linewidth}{@{}p{1.22in}X X@{}}
\toprule
\textbf{Method} & \textbf{Try when} & \textbf{Measure / watch}\\
\midrule
Statevector & Circuit is small enough for the available memory/time; useful as a reference on smaller examples. & Memory and runtime rise quickly with qubit count. Use it to validate smaller cases.\\
MPS & The state can be represented with manageable bond dimension; circuit structure and entanglement permit it. & Vary bond dimension only on controlled examples; compare with exact small answers or another reference.\\
Pauli-path & You need observable expectations and the circuit/observable structure is favorable for propagating terms. & Sweep truncation threshold. Smaller thresholds retain more terms and can cost more. Check per-circuit behavior.\\
Sampling & The task is about expectation values that can be estimated from measurement outcomes, or the platform provides sampled results. & Shot noise; basis changes for X/Y; distinguish sampling error from simulator approximation error.\\
\bottomrule
\end{tabularx}

Do not pick a backend from qubit count alone. Inspect the actual observable format, circuit gates/depth if available, connectivity if available, allowed backends, score, time limit, and budget. BlueQubit's docs currently expose statevector, MPS, and Pauli-path devices and an SDK call for cost/runtime estimates on some backends. Re-check live options and prices on challenge day.\footnote{\url{https://app.bluequbit.io/docs} and \url{https://app.bluequbit.io/sdk-docs/bluequbit.sdk.html}}

\section*{5. Team and agent workflow}
\textbf{Human roles (for a four-person team; adapt to actual event rules):}
\begin{itemize}
  \item \textbf{Captain / scorer:} reads the rules, owns the exact metric, makes the final method and spending calls, and submits.
  \item \textbf{Baseline engineer:} maintains the known-answer small-circuit estimator and qubit-order checks.
  \item \textbf{Simulator engineer:} tests BlueQubit methods and approximation controls on the same inputs.
  \item \textbf{Evaluator / recorder:} manages benchmark tables, job status, cost, confidence, and the submission explanation.
\end{itemize}

\textbf{Use coding agents as task workers, not as authorities.} Give each agent one bounded deliverable and an explicit acceptance check. Keep one human responsible for integrating each change. Example assignments:
\begin{enumerate}
  \item \textbf{Review the estimator:} ``Inspect this observable-estimation code for basis-rotation, Pauli-sign, and bit-order errors. Return the smallest example that demonstrates each suspected bug; do not rewrite unrelated code.''
  \item \textbf{Benchmark one method:} ``Using only the saved tutorial circuit and its listed Pauli sum, run the specified settings. Return a CSV with inputs, method settings, expectation, runtime, cost, and job ID. Do not launch any other jobs.''
  \item \textbf{Implement the scorer:} ``Implement the official scoring formula copied into this file. Add hand-calculated examples for each documented edge case. Do not infer missing rules.''
  \item \textbf{Independent challenge reader:} ``Read the active challenge rules. Extract circuit input format, observable input format, allowed methods, spending limits, deadline, output format, score formula, and any verification requirements. Link each claim to the exact page.''
\end{enumerate}

Store the BlueQubit token in an environment variable or local secrets file excluded from Git; do not paste it into the agent prompt, source code, shared notebook, or public logs. Set a human-approved maximum spend before any paid run. The BlueQubit SDK docs show examples for job submission, asynchronous wait/cancel, Pauli sums, MPS, Pauli-path, and the \texttt{estimate} call; use the active challenge's instructions for the actual permitted workflow.\footnote{\url{https://app.bluequbit.io/sdk-docs/index.html}}

\section*{6. Live benchmark worksheet}
Copy one row per run. Never compare results without preserving the inputs and settings.
\small
\begin{longtable}{@{}p{0.78in}p{0.62in}p{0.72in}p{0.77in}p{0.69in}p{0.68in}p{0.53in}p{0.58in}@{}}
\toprule
Circuit / obs. & Method & Setting & Estimate & Reference / score & Runtime & Cost & Job ID\\
\midrule
\endfirsthead
\toprule
Circuit / obs. & Method & Setting & Estimate & Reference / score & Runtime & Cost & Job ID\\
\midrule
\endhead
\rule{0pt}{0.55in} & & & & & & & \\
\rule{0pt}{0.55in} & & & & & & & \\
\rule{0pt}{0.55in} & & & & & & & \\
\rule{0pt}{0.55in} & & & & & & & \\
\rule{0pt}{0.55in} & & & & & & & \\
\rule{0pt}{0.55in} & & & & & & & \\
\bottomrule
\end{longtable}
\normalsize

\section*{7. Challenge-day first hour}
\begin{enumerate}
  \item Read the active rules and write down the exact metric, permitted backends, compute credits/spending rule, interface, deadline, result format, and verification requirements.
  \item Save a clean copy of the starter files and the first circuit/observable metadata. Do not hand-edit hidden inputs.
  \item Run one low-cost smoke test. Check input parsing, output dimensions/order, job completion, and score code.
  \item Produce a baseline for every required observable before spending time tuning one circuit.
  \item Select the next method from actual circuit metadata and your benchmark evidence. Change one setting at a time and log every run.
  \item Stop a method when it is consuming budget without improving the challenge metric or useful confidence. Keep enough time for reproducibility, explanation, and submission.
\end{enumerate}

\section*{8. Oral problem-solving drill}
Each person should answer these in under one minute, using the project code when helpful:
\begin{enumerate}
  \item In the team's own words, what does the challenge ask you to estimate?
  \item What does a Pauli string such as \texttt{XYZ} mean in the SDK input? Which qubit corresponds to each character in the actual API convention?
  \item Why is computational-basis measurement directly enough for Z but not for X or Y?
  \item How do you compute one Pauli-string sample from a measured bitstring?
  \item What evidence supports your chosen simulator and its settings?
  \item What changed in the score when you changed the setting, and what did it cost?
  \item Which output is an exact reference, which is approximate, and which is sampled?
  \item How can another teammate reproduce the result from the saved log?
\end{enumerate}

\section*{9. Do-not-skip checklist}
\begin{itemize}
  \item[\checkline] We copied the active scoring formula and implemented it literally.
  \item[\checkline] We ran the source's small known-answer examples and preserved outputs.
  \item[\checkline] We documented qubit indexing, bitstring order, Pauli-string order, and X/Y basis changes.
  \item[\checkline] We compared at least two methods on identical input where possible.
  \item[\checkline] We recorded runtime and cost/credits for every run and checked current limits.
  \item[\checkline] We bounded agent tasks and cloud spend; a human reviewed every merged result.
  \item[\checkline] Every teammate can explain the baseline, metric, selected method, and one limitation.
  \item[\checkline] We reserved time for submission checks and the explanation of how results were obtained.
\end{itemize}

\section*{10. Source shelf}
The exercises above are drawn from these linked first-party course pages and BlueQubit docs. Use the page's own exercise code/output as the answer check; versions and platform limits can change, so check the live pages again before running cloud jobs.
\begin{enumerate}[label={\textbf{S\arabic*.}},leftmargin=2.8em]
  \item \href{https://duquantum.org/}{DuQuantum 2026 official site}: event schedule, beginner track, event FAQ and team information.
  \item \href{https://qa.linkedin.com/company/bluequbit}{BlueQubit public announcement}: current public summary of the Duke challenge.
  \item \href{https://app.bluequbit.io/tutorials}{BlueQubit tutorials}: BQ 101, BQ 201, Pauli-path simulation and related hands-on examples.
  \item \href{https://app.bluequbit.io/sdk-docs/index.html}{BlueQubit Python SDK quickstart}: Qiskit circuits, MPS, Pauli-path observable run, async jobs.
  \item \href{https://app.bluequbit.io/sdk-docs/bluequbit.sdk.html}{BlueQubit SDK API reference}: device and truncation/bond-dimension options.
  \item \href{https://learning.quantum.ibm.com/course/utility-scale-quantum-computing/lesson-02-bits-gates-and-circuits}{IBM Quantum Learning, Lesson 2}: gate/statevector exercise and measurement basics.
  \item \href{https://learning.quantum.ibm.com/course/utility-scale-quantum-computing/lesson-03-teleportation}{IBM Quantum Learning, Lesson 3}: X/Y component tomography exercises.
  \item \href{https://learning.quantum.ibm.com/course/variational-algorithm-design/cost-functions}{IBM Quantum Learning, Cost Functions}: expectation values, Pauli decomposition, and a worked Hamiltonian example.
  \item \href{https://learning.quantum.ibm.com/tutorial/chsh-inequality}{IBM Quantum Learning, CHSH Inequality}: parameterized Bell-pair expectation-value sweep and post-processing.
  \item \href{https://quantum.cloud.ibm.com/docs/en/guides/specify-observables-pauli}{IBM Quantum documentation}: Pauli representation and basis changes for X, Y, and Z.
\end{enumerate}

\vfill
\begin{center}\small\textcolor{gray}{Prepared October 21, 2026. Public challenge details are time-sensitive; confirm the live rules at kickoff.}\end{center}

\end{document}
